\chapter{GRAVITATIONAL EQUILIBRIUM AND GRAVITATIONAL INSTABILITY}
\label{ch:13}

\section{Introduction}
\label{sec:116}

In this chapter we shall consider some general theorems on gravitational equilibrium and gravitational instability. These matters clearly have astronomical overtones; indeed, certain aspects of them have been extensively investigated in the astronomical contexts. However, our concern with them will be only to the extent they bear on the general theory of hydrodynamic and hydromagnetic stability.

\section{The virial theorem}
\label{sec:117}

Consider an inviscid fluid of zero electrical resistivity in which a magnetic field $\mathbf{H}(\mathbf{x})$ prevails. Suppose that the fluid is a perfect gas and that the ratio of the specific heats is $\gamma$. Suppose further that apart from the prevailing magnetic field and gas pressure the only other force acting on the medium is that derived from its own gravitation. Under these circumstances the equation of motion governing the fluid velocities is
\begin{equation}
\label{eq:13-1}
\rho \frac{du_i}{dt} = -\frac{\partial}{\partial x_k}\!\left(p + \frac{|\mathbf{H}|^2}{8\pi}\right) + \frac{\partial \mathfrak{B}}{\partial x_i} + \frac{1}{4\pi}\frac{\partial}{\partial x_j}H_i H_j,
\end{equation}
where
\begin{equation}
\label{eq:13-2}
\frac{d}{dt} = \frac{\partial}{\partial t} + u_j \frac{\partial}{\partial x_j}
\end{equation}
is the total time derivative. In equation~\eqref{eq:13-1} $\mathfrak{B}$ denotes the gravitational potential, and the rest of the symbols have their usual meanings. (We are setting $\mu = 1$ in the present analysis.)

\subsection*{(a) \textit{Some definitions and relations}}

Since we have supposed that the gravitational potential $\mathfrak{B}(\mathbf{x})$ is derived from the distribution of the matter present,
\begin{equation}
\label{eq:13-3}
\mathfrak{B}(\mathbf{x}) = G \int_V \frac{\rho(\mathbf{x}')}{|\mathbf{x}-\mathbf{x}'|}\,d\mathbf{x}',
\end{equation}
where for the sake of brevity we have written
\begin{equation}
\label{eq:13-4}
d\mathbf{x}' = dx_1\,dx_2\,dx_3
\end{equation}
and abridged three integral signs into one. In~\eqref{eq:13-3} the integration is effected over the entire volume $V$ occupied by the fluid.

We shall find it convenient to define the symmetric tensor,
\begin{equation}
\label{eq:13-5}
\mathfrak{B}_{ik}(\mathbf{x}) = G \int_V \rho(\mathbf{x}') \frac{(x_i-x_i')(x_k-x_k')}{|\mathbf{x}-\mathbf{x}'|^3}\,d\mathbf{x}',
\end{equation}
which represents a generalization of~\eqref{eq:13-3} and to which it reduces on contraction:
\begin{equation}
\label{eq:13-6}
\mathfrak{B}_{ii} = \mathfrak{B}.
\end{equation}
Similarly, in generalization of the usual definition of the gravitational potential energy, we shall define
\begin{equation}
\label{eq:13-7}
\mathfrak{W}_{ik} = -\tfrac{1}{2}G \iint_V \rho(\mathbf{x})\rho(\mathbf{x}') \frac{(x_i-x_i')(x_k-x_k')}{|\mathbf{x}-\mathbf{x}'|^3}\,d\mathbf{x}\,d\mathbf{x}';
\end{equation}
and the contraction of this tensor gives the gravitational potential energy
\begin{equation}
\label{eq:13-8}
\mathfrak{W}_{ii} = \mathfrak{W} = -\tfrac{1}{2}G \iint_V \frac{\rho(\mathbf{x})\rho(\mathbf{x}')}{|\mathbf{x}-\mathbf{x}'|}\,d\mathbf{x}\,d\mathbf{x}'.
\end{equation}
An identity which follows from the foregoing definitions is
\begin{equation}
\label{eq:13-9}
\mathfrak{W}_{ik} = \int_V \rho(\mathbf{x}) x_i \frac{\partial\mathfrak{B}}{\partial x_k}\,dx = \int_V \rho(\mathbf{x}) x_k \frac{\partial\mathfrak{B}}{\partial x_i}\,dx.
\end{equation}
This can be established as follows: by definition,
\begin{equation}
\label{eq:13-10}
\int dx\,\rho(\mathbf{x}) x_i \frac{\partial\mathfrak{B}}{\partial x_k} = G \int dx\,\rho(\mathbf{x}) x_i \frac{\partial}{\partial x_k} \int dx'\,\frac{\rho(\mathbf{x}')}{|\mathbf{x}-\mathbf{x}'|}
= -G \iint dx\,d\mathbf{x}'\,\rho(\mathbf{x})\rho(\mathbf{x}') \frac{x_i(x_k-x_k')}{|\mathbf{x}-\mathbf{x}'|^3}.
\end{equation}
By transposing the primed and the unprimed variables of integrations in~\eqref{eq:13-10} and taking the average of the two equivalent expressions, we verify that the double integral on the last line of~\eqref{eq:13-10} is the same as $\mathfrak{W}_{ik}$ defined in equation~\eqref{eq:13-7}.

In addition to $\mathfrak{B}_{ik}$ and $\mathfrak{W}_{ik}$, we shall find it useful to define the further tensors
\begin{equation}
\label{eq:13-11}
\mathfrak{T}_{ik} = \tfrac{1}{2}\int_V \rho u_i u_k\,d\mathbf{x}
\end{equation}
and
\begin{equation}
\label{eq:13-12}
\mathfrak{M}_{ik} = \frac{1}{8\pi}\int_V H_i H_k\,d\mathbf{x};\footnote{The specification of the volume $V$ over which this integration is effected requires some care; we return to this matter in \S\,(b) below.}
\end{equation}
and the contraction of these tensors gives the kinetic and the magnetic energies of the system
\begin{equation}
\label{eq:13-13}
\mathfrak{T} = \mathfrak{T}_{ii} = \tfrac{1}{2}\int_V \rho|\mathbf{u}|^2\,d\mathbf{x}
\end{equation}
and
\begin{equation}
\label{eq:13-14}
\mathfrak{M} = \mathfrak{M}_{ii} = \frac{1}{8\pi}\int_V |\mathbf{H}|^2\,d\mathbf{x}.
\end{equation}
The analogous expression for the internal heat energy of the system is
\begin{equation}
\label{eq:13-15}
\mathrm{U} = \frac{1}{\gamma-1}\int_V p\,d\mathbf{x}.
\end{equation}
Even as $\mathfrak{W}_{ik}$, $\mathfrak{T}_{ik}$, $\mathfrak{M}_{ik}$, and $\mathrm{U}$ characterize, in terms of a few parameters, the distribution of the different forms of energy in the configuration, so does the \textit{inertia tensor}
\begin{equation}
\label{eq:13-16}
I_{ik} = \int_V \rho x_i x_k\,d\mathbf{x}
\end{equation}
similarly characterize the distribution of the density in the configuration. The contraction of $I_{ik}$ gives the moment of inertia
\begin{equation}
\label{eq:13-17}
I = I_{ii} = \int_V \rho|\mathbf{x}|^2\,d\mathbf{x}.
\end{equation}

\subsection*{(b) \textit{The general form of the virial theorem}}

Returning to equation~\eqref{eq:13-1}, multiply it by $x_k$ and integrate it over the entire volume $V$ in which the fluid and the field pervade. The left-hand side of the equation can be reduced in the manner
\begin{equation}
\label{eq:13-18}
\int_V \rho x_k \frac{du_i}{dt}\,d\mathbf{x} = \int_V \rho x_k \frac{d^2 x_i}{dt^2}\,d\mathbf{x} = \frac{d}{dt}\!\left(\int_V \rho \frac{d}{dt}(x_i x_k)\right) - \int_V \rho \frac{dx_i}{dt}\frac{dx_k}{dt}\,d\mathbf{x}
= \int_V \rho \frac{d}{dt}\!\left(x_k\frac{dx_i}{dt}\right)d\mathbf{x} - 2\mathfrak{T}_{ik}.
\end{equation}
The terms on the right-hand side, similarly, give
\begin{align}
\label{eq:13-19}
-\int_V x_k \frac{\partial}{\partial x_i}\!\left(p+\frac{|\mathbf{H}|^2}{8\pi}\right)d\mathbf{x}
&= -\int_S\!\left(p+\frac{|\mathbf{H}|^2}{8\pi}\right)x_k\,dS_i + \delta_{ik}\int_V\!\left(p+\frac{|\mathbf{H}|^2}{8\pi}\right)d\mathbf{x} \nonumber\\
&= -\int_S\!\left(p+\frac{|\mathbf{H}|^2}{8\pi}\right)x_k\,dS_i + \delta_{ik}((\gamma-1)\mathrm{U}+\mathfrak{M}),
\end{align}
\begin{equation}
\label{eq:13-20}
\int_V \rho x_k \frac{\partial\mathfrak{B}}{\partial x_i}\,d\mathbf{x} = \mathfrak{W}_{ik},
\end{equation}
and
\begin{equation}
\label{eq:13-21}
\frac{1}{4\pi}\int_V x_k \frac{\partial}{\partial x_j}H_i H_j\,d\mathbf{x} = \frac{1}{4\pi}\int_S x_k H_i H_j\,dS_j - \frac{1}{4\pi}\int_V H_i H_k\,d\mathbf{x}
= \frac{1}{4\pi}\int_S x_k H_i H_j\,dS_j - 2\mathfrak{M}_{ik},
\end{equation}
where in equations~\eqref{eq:13-19} and~\eqref{eq:13-21} the surface integrals are extended over the surface $S$ bounding $V$.

Combining equations~\eqref{eq:13-18}--\eqref{eq:13-21}, we obtain
\begin{equation}
\label{eq:13-22}
\int_V \rho \frac{d}{dt}\!\left(x_k\frac{dx_i}{dt}\right)d\mathbf{x}
= 2\mathfrak{T}_{ik} + \mathfrak{W}_{ik} - 2\mathfrak{M}_{ik} + \delta_{ik}((\gamma-1)\mathrm{U}+\mathfrak{M}) + \frac{1}{4\pi}\int_S x_k H_i H_j\,dS_j - \int_S px_k\,dS_i.
\end{equation}

We shall suppose that the volume $V$, over which the integrations are extended, includes the whole system so that all the variables, $\rho$, $p$, and $\mathbf{H}$ may be assumed to vanish on $S$. It is important to note that this definition of $V$ may require us to include in it volumes which one may normally consider as external to the `natural' boundary of the system, namely, the surface on which the density $\rho$ and the material pressure $p$ vanish. The assumption that $\mathbf{H}$ vanishes on $S$ may, indeed, require us to place $S$ at infinity. This latter possibility arises because magnetic fields can extend far beyond the conventional limits of a material object; but to the extent the object is the seat of the magnetic field, there is justification in including all portions of space into which the field extends as parts of the system. And since the field of an object isolated in space must decrease at least as rapidly as that of a dipole (i.e., as $r^{-3}$), the surface integrals over the components of $\mathbf{H}$ in equation~\eqref{eq:13-22} will vanish under these circumstances when $S \to \infty$. However, it may sometimes be convenient to let $S$ coincide with the natural boundary, in which case the surface integrals over $\mathbf{H}$ in equation~\eqref{eq:13-22} must be retained; and, moreover, $\mathfrak{M}_{ik}$ will then refer to only that part of the field which is interior to $S$.

In the present analysis we shall suppose that $V$ includes (as we have already remarked) all parts of space in which the fluid and the field pervade. With this explicit understanding, equation~\eqref{eq:13-22} becomes
\begin{equation}
\label{eq:13-23}
\int_V \frac{d}{dt}\!\left(x_k\frac{dx_i}{dt}\right)d\mathbf{x}
= 2\mathfrak{T}_{ik} - 2\mathfrak{M}_{ik} + \mathfrak{W}_{ik} + \delta_{ik}((\gamma-1)\mathrm{U}+\mathfrak{M}).
\end{equation}

Since all the tensors on the right-hand side of this equation are symmetric, the tensor on the left-hand side of the equation must also be symmetric. Therefore,
\begin{equation}
\label{eq:13-24}
\int_V \rho \frac{d}{dt}\!\left(x_k\frac{dx_i}{dt}\right)d\mathbf{x} = \int_V \rho \frac{d}{dt}\!\left(x_i\frac{dx_k}{dt}\right)d\mathbf{x}.
\end{equation}
An immediate consequence of this result is
\begin{equation}
\label{eq:13-25}
\int_V \rho \frac{d}{dt}\!\left(x_i\frac{dx_k}{dt} - x_k\frac{dx_i}{dt}\right)d\mathbf{x} = \frac{d}{dt}\int_V \rho\!\left(x_i\frac{dx_k}{dt} - x_k\frac{dx_i}{dt}\right)d\mathbf{x} = 0,
\end{equation}
where in taking $d/dt$ outside of the integral sign, we have made use of the constancy of $\rho\,d\mathbf{x}$ assured by the equation of continuity. \textit{Equation~\eqref{eq:13-25} expresses the constancy of the total angular momentum of the system.} It is worth noting that the existence of this integral of the equations of motion has not been affected by the presence of the magnetic field.

A further consequence of the identity~\eqref{eq:13-24} is that the quantity on the left-hand side of equation~\eqref{eq:13-23} can be replaced by
\begin{equation}
\label{eq:13-26}
\tfrac{1}{2}\int_V \rho \frac{d^2}{dt^2}(x_i x_k)\,d\mathbf{x} = \tfrac{1}{2}\frac{d^2}{dt^2}\int_V \rho x_i x_k\,d\mathbf{x} = \tfrac{1}{2}\frac{d^2 I_{ik}}{dt^2}.
\end{equation}
With this replacement, equation~\eqref{eq:13-23} gives
\begin{equation}
\label{eq:13-27}
\frac{1}{2}\frac{d^2 I_{ik}}{dt^2} = 2\mathfrak{T}_{ik} - 2\mathfrak{M}_{ik} + \mathfrak{W}_{ik} + \delta_{ik}((\gamma-1)\mathrm{U}+\mathfrak{M}).
\end{equation}
This equation represents the complete statement of \textit{the virial theorem for hydromagnetics}.

When $i \neq k$, equation~\eqref{eq:13-27} gives
\begin{equation}
\label{eq:13-28}
\frac{1}{2}\frac{d^2 I_{ik}}{dt^2} = 2\mathfrak{T}_{ik} - 2\mathfrak{M}_{ik} + \mathfrak{W}_{ik} \quad (i \neq k);
\end{equation}
and by contracting the indices in equation~\eqref{eq:13-27}, we obtain
\begin{equation}
\label{eq:13-29}
\frac{1}{2}\frac{d^2 I}{dt^2} = 2\mathfrak{T}+\mathfrak{M}+\mathfrak{W}+3(\gamma-1)\mathrm{U}.
\end{equation}

\subsection*{(c) \textit{The virial theorem for equilibrium configurations}}

For configurations in equilibrium and in a steady state, equations~\eqref{eq:13-27} and~\eqref{eq:13-29} give
\begin{equation}
\label{eq:13-30}
2\mathfrak{T}_{ik} - 2\mathfrak{M}_{ik} + \mathfrak{W}_{ik} + \delta_{ik}((\gamma-1)\mathrm{U}+\mathfrak{M}) = 0
\end{equation}
and
\begin{equation}
\label{eq:13-31}
2\mathfrak{T}+\mathfrak{M}+\mathfrak{W}+3(\gamma-1)\mathrm{U} = 0.
\end{equation}
An alternative way of stating the result expressed by equation~\eqref{eq:13-30} is that \textit{the tensor} $2(\mathfrak{T}_{ik}-\mathfrak{M}_{ik})+\mathfrak{W}_{ik}$ \textit{is isotropic}.

Consider first the case when there are no fluid motions or magnetic fields. In this case, equation~\eqref{eq:13-30} gives
\begin{equation}
\label{eq:13-32}
\mathfrak{W}_{ik} = \int_V \rho(\mathbf{x}) x_i \frac{\partial\mathfrak{B}}{\partial x_k}\,d\mathbf{x} = -(\gamma-1)\mathrm{U}\delta_{ik}.
\end{equation}
This relation is compatible with a spherical symmetry of the underlying density distribution but does not require it.

When $\mathfrak{T}_{ik}$ and $\mathfrak{M}_{ik}$ are not zero, $\mathfrak{W}_{ik}$ will not, in general, be diagonal. Consequently, \textit{a spherical symmetry of the configuration is, in general, incompatible with the presence of fluid motions and magnetic fields.} An exception is possible if $\mathfrak{T}_{ik} = \mathfrak{M}_{ik}$---an identity which is satisfied, for example, by the equipartition solution discussed in Chapter~XII, \S\,113; for, when $\mathfrak{T}_{ik} = \mathfrak{M}_{ik}$, equation~\eqref{eq:13-30} gives
\begin{equation}
\label{eq:13-33}
\mathfrak{W}_{ik} = \int_V \rho(\mathbf{x}) x_i \frac{\partial\mathfrak{B}}{\partial x_k}\,d\mathbf{x} = -\{(\gamma-1)\mathrm{U}+\mathfrak{M}\}\delta_{ik};
\end{equation}
and this equation is not incompatible with spherical symmetry.

In the special case when only magnetic fields are present,
\begin{equation}
\label{eq:13-34}
\mathfrak{W}_{ik} = 2\mathfrak{M}_{ik} - \delta_{ik}((\gamma-1)\mathrm{U}+\mathfrak{M});
\end{equation}
and the impossibility of spherical symmetry, in general, follows from the relation
\begin{equation}
\label{eq:13-35}
\mathfrak{W}_{ik} = 2\mathfrak{M}_{ik} \neq 0 \quad \text{for } i \neq k.
\end{equation}

Consider next the contracted form~\eqref{eq:13-31} of the general relation. If $\mathfrak{E}$ denotes the total energy of the configuration,
\begin{equation}
\label{eq:13-36}
\mathfrak{E} = \mathfrak{T}+\mathrm{U}+\mathfrak{M}+\mathfrak{W}.
\end{equation}
From equations~\eqref{eq:13-31} and~\eqref{eq:13-36} it would appear that, as far as these scalar equations go, the magnetic energy can be considered together with the gravitational potential energy. Thus, by eliminating $\mathfrak{M}+\mathfrak{W}$ between the two equations, we obtain the relation,
\begin{equation}
\label{eq:13-37}
\mathfrak{E} = -\mathfrak{T}-(3\gamma-4)\mathrm{U},
\end{equation}
which is independent of $\mathfrak{M}$. Alternatively, by eliminating $\mathrm{U}$, we obtain
\begin{equation}
\label{eq:13-38}
\mathfrak{E} = \frac{1}{3(\gamma-1)}\{(3\gamma-4)(\mathfrak{M}+\mathfrak{W}) + (3\gamma-5)\mathfrak{T}\}.
\end{equation}

In case $\mathfrak{T} = 0$, equations~\eqref{eq:13-37} and~\eqref{eq:13-38} give
\begin{equation}
\label{eq:13-39}
\mathfrak{E} = -(3\gamma-4)\mathrm{U}
\end{equation}
and
\begin{equation}
\label{eq:13-40}
\mathfrak{E} = \frac{3\gamma-4}{3(\gamma-1)}(\mathfrak{M}+\mathfrak{W}).
\end{equation}

An immediate consequence of equation~\eqref{eq:13-39} is this: in a configuration for which $\gamma = \frac{4}{3}$, $\mathfrak{E} = 0$ in a steady state; and this is true independently of the radius of the configuration. A small radial expansion during which the configuration changes from one equilibrium state to an adjacent equilibrium state without any expenditure of energy is, therefore, possible. Moreover, since a gravitationally stable configuration must have a total energy which is negative, it is clearly necessary that \textit{for stability the ratio of the specific heats $\gamma$ must be greater than $\frac{4}{3}$}.

A further consequence of equation~\eqref{eq:13-40} is that for the total energy $\mathfrak{E}$ to be negative it is necessary that
\begin{equation}
\label{eq:13-41}
\mathfrak{M}+\mathfrak{W} < 0 \quad \text{or} \quad \mathfrak{M} < |\mathfrak{W}|.
\end{equation}
Therefore, \textit{for the existence of a stable equilibrium, it is necessary that the total magnetic energy of a system does not exceed its negative gravitational potential energy.}

The inequality~\eqref{eq:13-41} enables one to place a useful upper limit on the possible strengths of the magnetic fields of cosmical bodies. Thus, by writing
\begin{equation}
\label{eq:13-42}
\mathfrak{M} = \tfrac{1}{6}\bar{H}^2\rangle_{av} R^3 \quad \text{and} \quad |\mathfrak{W}| = q\,\frac{GM^2}{R},
\end{equation}
where $R$ is a measure of the linear dimension of the system, $M$ is its mass, and $q$ is a numerical constant of order unity, we obtain
\begin{equation}
\label{eq:13-43}
\sqrt{\langle H^2\rangle_{av}} < \frac{M}{R^2}\sqrt{6qG}.
\end{equation}

\section{The virial theorem for small oscillations about equilibrium}
\label{sec:118}

We shall now consider small oscillations about equilibrium of a configuration in which magnetic fields are present. We shall, however, suppose that in the stationary state there are no fluid motions; and we shall seek the form which the virial theorem takes under these circumstances.

Considering periodic oscillations with a gyration frequency $\sigma$, let $\boldsymbol{\xi}e^{i\sigma t}$ denote the Lagrangian displacement of an element of mass, $dm = \rho\,d\mathbf{x}$, from its equilibrium position at $\mathbf{x}$. Let $\delta\rho e^{i\sigma t}$, $\delta\rho e^{i\sigma t}$, and $\delta\mathbf{H}e^{i\sigma t}$ denote the corresponding changes in the other physical variables as we follow the fluid element with its motion. The equation of continuity ensures the constancy of $dm$ for such Lagrangian displacements; and if we suppose that the oscillations take place adiabatically, then
\begin{equation}
\label{eq:13-44}
\delta\rho/\rho = \gamma\,\delta\rho/\rho,
\end{equation}
where
\begin{equation}
\label{eq:13-45}
\delta\rho/\rho = -\operatorname{div}\boldsymbol{\xi}.
\end{equation}

If $\delta I_{ik}\,e^{i\sigma t}$, $\delta\mathfrak{W}_{ik}\,e^{i\sigma t}$, etc., denote the first-order changes in the various integrals representing these quantities in the stationary state, then equation~\eqref{eq:13-27} gives
\begin{equation}
\label{eq:13-46}
-\tfrac{1}{2}\sigma^2 \delta I_{ik} = \delta\mathfrak{W}_{ik} + \delta_{ik}((\gamma-1)\delta\mathrm{U}+\delta\mathfrak{M}) - 2\delta\mathfrak{M}_{ik};
\end{equation}
the term in $\mathfrak{T}_{ik}$ makes no contribution in this order since we have supposed that there are no zero-order fluid motions.

We shall now consider, in turn, the various quantities occurring in equation~\eqref{eq:13-46}. Clearly,
\begin{equation}
\label{eq:13-47}
\delta I_{ik} = \int_V \rho(\xi_i x_k + \xi_k x_i)\,d\mathbf{x}.
\end{equation}
Also, according to the definitions of $\mathfrak{B}_{ik}$ and $\mathfrak{W}_{ik}$ given in equations~\eqref{eq:13-5} and~\eqref{eq:13-7},
\begin{equation}
\label{eq:13-48}
\delta\mathfrak{W}_{ik} = -G\iint_V dm\,dm'\,\xi_j\frac{\partial}{\partial x_j}\frac{(x_i-x_i')(x_k-x_k')}{|\mathbf{x}-\mathbf{x}'|^3}
= -G\int_V d\mathbf{x}\,\rho(\mathbf{x})\xi_j\frac{\partial}{\partial x_j}\int_V d\mathbf{x}'\,\rho(\mathbf{x}')\frac{(x_i-x_i')(x_k-x_k')}{|\mathbf{x}-\mathbf{x}'|^3}
= -\int_V d\mathbf{x}\,\rho(\mathbf{x})\xi_j\frac{\partial\mathfrak{B}_{ik}}{\partial x_j}.
\end{equation}
This last relation is a generalization of the familiar formula,
\begin{equation}
\label{eq:13-49}
\delta\mathfrak{W} = -\int_V d\mathbf{x}\,\rho(\mathbf{x})\xi_j\frac{\partial\mathfrak{B}}{\partial x_j},
\end{equation}
for the first-order change in the gravitational potential energy consequent to a slight redistribution of the matter in the system.

Turning next to the change in the internal energy, we have
\begin{equation}
\label{eq:13-50}
(\gamma-1)\delta\mathrm{U} = \int_V \frac{\delta p}{\rho}\,\rho\,d\mathbf{x},
\end{equation}
or, making use of the relations~\eqref{eq:13-44} and~\eqref{eq:13-45}, we have
\begin{equation}
\label{eq:13-51}
(\gamma-1)\delta\mathrm{U} = (\gamma-1)\int_V \frac{\gamma\,\delta\rho}{\rho}\,p\,d\mathbf{x} = -(\gamma-1)\int_V p\frac{\partial\xi_j}{\partial x_j}\,d\mathbf{x}.
\end{equation}
Letting
\begin{equation}
\label{eq:13-52}
\Pi = p + |\mathbf{H}|^2/8\pi,
\end{equation}
we can write
\begin{equation}
\label{eq:13-53}
\delta\mathrm{U} = -\int_V \Pi\frac{\partial\xi_j}{\partial x_j}\,d\mathbf{x} + \frac{1}{8\pi}\int_V |\mathbf{H}|^2\frac{\partial\xi_j}{\partial x_j}\,d\mathbf{x}.
\end{equation}

Integrating by parts the first of the two integrals on the right-hand side, we obtain
\begin{equation}
\label{eq:13-54}
\int_V \Pi\frac{\partial\xi_j}{\partial x_j}\,d\mathbf{x} = \int_S \Pi\,\xi_j\,dS_j - \int_V \xi_j\frac{\partial\Pi}{\partial x_j}\,d\mathbf{x}.
\end{equation}
We shall suppose that $S$ is so placed that all the physical variables and their perturbations vanish on $S$. (As we have remarked in \S\,117, this may require that $S$ be placed at infinity.) Then
\begin{equation}
\label{eq:13-55}
\int_V \Pi\frac{\partial\xi_j}{\partial x_j}\,d\mathbf{x} = -\int_V \xi_j\frac{\partial\Pi}{\partial x_j}\,d\mathbf{x},
\end{equation}
or, making use of the relation,
\begin{equation}
\label{eq:13-56}
\frac{\partial\Pi}{\partial x_j} = \rho\frac{\partial\mathfrak{B}}{\partial x_j} + \frac{1}{4\pi}\frac{\partial}{\partial x_j}H_i H_j,
\end{equation}
which obtains in equilibrium, we have
\begin{equation}
\label{eq:13-57}
\int_V \Pi\frac{\partial\xi_j}{\partial x_j}\,d\mathbf{x} = -\int_V \rho\xi_j\frac{\partial\mathfrak{B}}{\partial x_j}\,d\mathbf{x} - \frac{1}{4\pi}\int_V \xi_j\frac{\partial}{\partial x_i}H_i H_j\,d\mathbf{x}.
\end{equation}
Integrating by parts the second of the two integrals on the right-hand side of this equation, we obtain
\begin{equation}
\label{eq:13-58}
\int_V \Pi\frac{\partial\xi_j}{\partial x_j}\,d\mathbf{x} = -\int_V \rho\xi_j\frac{\partial\mathfrak{B}}{\partial x_j}\,d\mathbf{x} + \frac{1}{4\pi}\int_V H_i H_j\frac{\partial\xi_j}{\partial x_i}\,d\mathbf{x};
\end{equation}
the integrated part, again, making no contribution. Now combining equations~\eqref{eq:13-53} and~\eqref{eq:13-58}, we have
\begin{equation}
\label{eq:13-59}
\delta\mathrm{U} = \int_V \rho\xi_j\frac{\partial\mathfrak{B}}{\partial x_j}\,d\mathbf{x} + \frac{1}{8\pi}\int_V\!\left(|\mathbf{H}|^2\frac{\partial\xi_j}{\partial x_j} - 2H_i H_j\frac{\partial\xi_j}{\partial x_i}\right)d\mathbf{x}.
\end{equation}

Finally, considering $\delta\mathfrak{M}_{ik}$, we have
\begin{equation}
\label{eq:13-60}
\delta\mathfrak{M}_{ik} = \frac{1}{8\pi}\int_V (H_i\,\delta H_k + H_k\,\delta H_i)\,d\mathbf{x} + \frac{1}{8\pi}\int_V H_i H_k\frac{\partial\xi_j}{\partial x_j}\,d\mathbf{x},
\end{equation}
where the second term arises from allowing for the variation of the volume element $d\mathbf{x}$, following the motion in accordance with equation~\eqref{eq:13-45}, and the constancy of $dm = \rho\,d\mathbf{x}$.

Now the change in the magnetic field $\delta\mathbf{H}$, as we follow the motion of the fluid element, is given by
\begin{equation}
\label{eq:13-61}
\delta\mathbf{H} = \operatorname{curl}(\boldsymbol{\xi}\times\mathbf{H}) + (\boldsymbol{\xi}\cdot\operatorname{grad})\mathbf{H}.
\end{equation}
The first term on the right-hand side gives the Eulerian change in $\mathbf{H}$ at a \textit{fixed point} caused by the redistribution of the matter, while the second term is the allowance for the fact that, as we follow the fluid element, it finds itself in a slightly displaced location. In the notation of Cartesian tensors, equation~\eqref{eq:13-61} becomes
\begin{equation}
\label{eq:13-62}
\delta H_i = H_j\frac{\partial\xi_i}{\partial x_j} - H_i\frac{\partial\xi_j}{\partial x_j} + \xi_j\frac{\partial H_i}{\partial x_j}.
\end{equation}
Making use of this relation, we find that equation~\eqref{eq:13-60} reduces to
\begin{equation}
\label{eq:13-63}
\delta\mathfrak{M}_{ik} = \frac{1}{8\pi}\int_V\!\left(H_j H_k\frac{\partial\xi_i}{\partial x_j} + H_i H_j\frac{\partial\xi_k}{\partial x_j} - H_i H_k\frac{\partial\xi_j}{\partial x_j}\right)d\mathbf{x};
\end{equation}
and by contracting, we find
\begin{equation}
\label{eq:13-64}
\delta\mathfrak{M} = -\frac{1}{8\pi}\int_V\!\left(|\mathbf{H}|^2\frac{\partial\xi_j}{\partial x_j} - 2H_i H_j\frac{\partial\xi_i}{\partial x_j}\right)d\mathbf{x}.
\end{equation}

Now inserting in equation~\eqref{eq:13-46} for $\delta I_{ik}$, $\delta\mathfrak{W}_{ik}$, etc., the expressions we have derived for them, we find
\begin{equation}
\label{eq:13-65}
\begin{split}
-\tfrac{1}{2}\sigma^2\int_V \rho(\xi_i x_k + \xi_k x_i)\,d\mathbf{x}
&= -\int_V \rho\xi_j\frac{\partial\mathfrak{B}_{ik}}{\partial x_j}\,d\mathbf{x} + \delta_{ik}(\gamma-1)\int_V \rho\xi_j\frac{\partial\mathfrak{B}}{\partial x_j}\,d\mathbf{x} + \\
&\quad + \delta_{ik}\frac{\gamma-2}{8\pi}\int_V\!\left(|\mathbf{H}|^2\frac{\partial\xi_j}{\partial x_j} - 2H_i H_j\frac{\partial\xi_j}{\partial x_i}\right)d\mathbf{x} + \\
&\quad + \frac{1}{4\pi}\int_V\!\left(H_i H_j\frac{\partial\xi_k}{\partial x_j} + H_k H_j\frac{\partial\xi_i}{\partial x_j} - H_i H_k\frac{\partial\xi_j}{\partial x_j}\right)d\mathbf{x}.
\end{split}
\end{equation}
By contracting this equation, we obtain
\begin{equation}
\label{eq:13-66}
-\sigma^2\int_V \rho\xi_i x_i\,d\mathbf{x}
= (3\gamma-4)\int_V \rho\xi_j\frac{\partial\mathfrak{B}}{\partial x_j}\,d\mathbf{x} + \frac{1}{8\pi}\int_V\!\left(|\mathbf{H}|^2\frac{\partial\xi_j}{\partial x_j} - 2H_i H_j\frac{\partial\xi_i}{\partial x_j}\right)d\mathbf{x}.
\end{equation}

\subsection*{(a) \textit{A characteristic equation for determining the periods of oscillation}}

Equations~\eqref{eq:13-65} and~\eqref{eq:13-66} can be used to obtain estimates for $\sigma^2$ by inserting in them suitable `trial' functions for $\boldsymbol{\xi}$. Thus, the simplest assumption
\begin{equation}
\label{eq:13-67}
\boldsymbol{\xi} = \text{constant} \times \mathbf{x}
\end{equation}
together with equation~\eqref{eq:13-66}, yields the formula
\begin{equation}
\label{eq:13-68}
\sigma^2 = -(3\gamma-4)\frac{\mathfrak{W}+\mathfrak{M}}{I}.
\end{equation}
However, since a configuration with a prevalent magnetic field cannot be spherically symmetric, the use of a trial function which corresponds to a uniform radial expansion would appear unsuitable; indeed, the substitution of~\eqref{eq:13-67} in the tensor equation~\eqref{eq:13-65} will lead to gross inconsistencies unless the departures from spherical symmetry of the equilibrium are small. If the latter should not be the case, a more consistent procedure would appear to be the following.

Assume that
\begin{equation}
\label{eq:13-69}
\boldsymbol{\xi} = \mathsf{X}\mathbf{x},
\end{equation}
where $\mathsf{X}$ is a symmetric matrix:
\begin{equation}
\label{eq:13-70}
X_{ij} = X_{ji}.
\end{equation}
Further, define the super-matrix
\begin{equation}
\label{eq:13-71}
\mathfrak{W}_{ij;ik} = \int_V \rho(\mathbf{x}) x_j \frac{\partial^2\mathfrak{B}_{ik}}{\partial x_j \partial x_l}\,d\mathbf{x}.
\end{equation}
By contracting this super-matrix with respect to $i$ and $k$, we obtain
\begin{equation}
\label{eq:13-72}
\mathfrak{W}_{ij;il} = \int_V \rho(\mathbf{x}) x_j \frac{\partial^2\mathfrak{B}}{\partial x_j \partial x_l}\,d\mathbf{x} = \mathfrak{W}_{jl}.
\end{equation}
Also, let
\begin{equation}
\label{eq:13-73}
\boldsymbol{\mathfrak{M}} = (\mathfrak{M}_{ik}), \quad \mathbf{W} = (\mathfrak{W}_{ik}), \quad \text{and} \quad \mathbf{I} = (I_{ik}).
\end{equation}
Clearly,
\begin{equation}
\label{eq:13-74}
\mathfrak{M} = \mathfrak{M}_{ii} = \operatorname{tr}(\boldsymbol{\mathfrak{M}}) \quad \text{and} \quad \mathfrak{W} = \mathfrak{W}_{ii} = \operatorname{tr}(\mathbf{W}),
\end{equation}
where `tr' stands for trace (meaning diagonal sum). Similarly,
\begin{equation}
\label{eq:13-75}
\frac{\partial\xi_i}{\partial x_j} = \operatorname{tr}(\mathsf{X}).
\end{equation}

With the foregoing definitions, the substitution of the trial function~\eqref{eq:13-69} in equation~\eqref{eq:13-65} leads to the result
\begin{equation}
\label{eq:13-76}
\begin{split}
-\tfrac{1}{2}\sigma^2(\mathbf{I}\mathsf{X}+\mathsf{X}\mathbf{I})_{ik}
&= -X_{jl}\,\mathfrak{W}_{ij;lk} + \\
&\quad + \delta_{ik}\{(\gamma-1)\operatorname{tr}(\mathsf{X}\mathbf{W}) + (\gamma-2)[\operatorname{tr}(\mathsf{X})\operatorname{tr}(\boldsymbol{\mathfrak{M}}) - 2\operatorname{tr}(\mathsf{X}\boldsymbol{\mathfrak{M}})]\} + \\
&\quad + 2\{\operatorname{tr}(\mathsf{X})\mathfrak{M}_{ik} - [\mathsf{X}\boldsymbol{\mathfrak{M}} + \boldsymbol{\mathfrak{M}}\mathsf{X}]_{ik}\}.
\end{split}
\end{equation}
Equation~\eqref{eq:13-76} represents a system of linear homogeneous equations for the six coefficients of the (assumed) symmetric linear transformation $\mathsf{X}$. The determinant of the system must, therefore, vanish; and the resulting characteristic equation will not only determine $\sigma^2$ but also the transformation $\mathsf{X}$ (apart from a constant of proportionality).

Finally, we may note that by contracting the tensor equation~\eqref{eq:13-76} we obtain the scalar equation
\begin{equation}
\label{eq:13-77}
-\sigma^2 \operatorname{tr}(\mathbf{I}\mathsf{X}) = (3\gamma-4)\{\operatorname{tr}(\mathsf{X}\mathbf{W})+\operatorname{tr}(\mathsf{X})\operatorname{tr}(\boldsymbol{\mathfrak{M}}) - 2\operatorname{tr}(\mathsf{X}\boldsymbol{\mathfrak{M}})\}
\end{equation}
which is a generalization of~\eqref{eq:13-68}.

\section{The gravitational instability of an infinite homogeneous medium; Jeans's criterion}
\label{sec:119}

A particularly simple example of gravitational instability was discovered by Jeans. In this example, we start from an infinite homogeneous medium at rest and consider the velocity of propagation of a small fluctuation in the density. If the gravitational effects of the fluctuation are ignored, the problem reduces to the classical one of the propagation of sound; and as is well known, the velocity of sound $c$ is independent of the wave number and is given by
\begin{equation}
\label{eq:13-78}
c^2 = \gamma p/\rho.
\end{equation}
If the change in the gravitational potential consequent to the fluctuation $\delta\rho$ in the density is taken into account, the velocity of wave propagation depends, as we shall presently see, on the wave number $k$, and is, indeed, imaginary for all wave numbers less than a certain value $k_J$. The instability which follows for $k < k_J$ is the gravitational instability discovered by Jeans.

The demonstration of the gravitational instability of an infinite homogeneous medium is simple. The relevant perturbation equations are
\begin{equation}
\label{eq:13-79}
\rho\frac{\partial\mathbf{u}}{\partial t} = -\operatorname{grad}\delta p + \rho\operatorname{grad}\delta\mathfrak{B},
\end{equation}
\begin{equation}
\label{eq:13-80}
\frac{\partial}{\partial t}\delta\rho = -\rho\operatorname{div}\mathbf{u},
\end{equation}
and
\begin{equation}
\label{eq:13-81}
\nabla^2\delta\mathfrak{B} = -4\pi G\delta\rho.
\end{equation}
Further, we shall suppose that the fluctuations in the pressure and the density take place adiabatically so that
\begin{equation}
\label{eq:13-82}
\delta p/p = \gamma\,\delta\rho/\rho
\end{equation}
or
\begin{equation}
\label{eq:13-83}
\delta p = c^2\delta\rho.
\end{equation}
Inserting this last expression for $\delta p$ in equation~\eqref{eq:13-79} and taking the divergence of the resulting equation, we obtain
\begin{equation}
\label{eq:13-84}
\rho\frac{\partial}{\partial t}\operatorname{div}\mathbf{u} = -c^2\nabla^2\delta\rho + \rho\nabla^2\delta\mathfrak{B}.
\end{equation}
Now making use of equations~\eqref{eq:13-80} and~\eqref{eq:13-81}, we find
\begin{equation}
\label{eq:13-85}
\frac{\partial^2}{\partial t^2}\delta\rho = c^2\nabla^2\delta\rho + 4\pi G\rho\,\delta\rho.
\end{equation}
This equation allows solutions of the form
\begin{equation}
\label{eq:13-86}
\delta\rho = \text{constant}\; e^{i(\mathbf{k}\cdot\mathbf{x}+\sigma t)}
\end{equation}
provided
\begin{equation}
\label{eq:13-87}
\sigma^2 = c^2 k^2 - 4\pi G\rho.
\end{equation}

The velocity of propagation $V_p$ of these waves is, therefore, given by
\begin{equation}
\label{eq:13-88}
V_p = \frac{\sigma}{k} = c\sqrt{1-\frac{4\pi G\rho}{c^2 k^2}}.
\end{equation}
We may call $V_J$ the \textit{Jeans velocity}. From equation~\eqref{eq:13-87} or~\eqref{eq:13-88} it follows that \textit{there is instability for all wave numbers}
\begin{equation}
\label{eq:13-89}
k < k_J \quad \text{where} \quad k_J = \tfrac{1}{c}\sqrt{4\pi G\rho}.
\end{equation}
This is Jeans's result; and the condition~\eqref{eq:13-89} for instability is often referred to as \textit{Jeans's criterion}.

The origin of the gravitational instability associated with the long waves lies in the circumstance that, while the generation of a pressure wave requires the expenditure of energy proportional to $c^2$, it simultaneously results in a release of gravitational potential energy of an amount proportional to $4\pi G\rho/k^2$; and when Jeans's criterion is satisfied, the system can spontaneously go to states of lower energy with the liberation of thermal energy; and this, of course, means instability.

\section{The effect of a uniform rotation and a uniform magnetic field on Jeans's criterion}
\label{sec:120}

We shall prove that \textit{Jeans's criterion for the gravitational instability of an infinite homogeneous medium is unaffected by the presence, separately, or simultaneously, of a uniform rotation and a uniform magnetic field.}

\subsection*{(a) \textit{The effect of a uniform rotation}}

The effect of a uniform rotation $\boldsymbol{\Omega}$ on Jeans's criterion is readily determined. In a frame of reference rotating with an angular velocity $\boldsymbol{\Omega}$, equation~\eqref{eq:13-79} is replaced by
\begin{equation}
\label{eq:13-90}
\rho\frac{\partial\mathbf{u}}{\partial t} = -\operatorname{grad}\delta p + \rho\operatorname{grad}\delta\mathfrak{B} + 2\rho\,\mathbf{u}\times\boldsymbol{\Omega}
\end{equation}
and equations~\eqref{eq:13-80} and~\eqref{eq:13-81} are unaffected.

We shall consider solutions of equations~\eqref{eq:13-80},~\eqref{eq:13-81}, and~\eqref{eq:13-90} which correspond to the propagation of waves in the $z$-direction (say). For such solutions $\partial/\partial z$ is the only non-vanishing component of the gradient; and in a system of coordinates so oriented that
\begin{equation}
\label{eq:13-91}
\boldsymbol{\Omega} = (0, \Omega_y, \Omega_z),
\end{equation}
the perturbation equations become
\begin{equation}
\label{eq:13-92}
\frac{\partial u_x}{\partial t} + 2\Omega_y u_z - 2\Omega_z u_y = 0,
\end{equation}
\begin{equation}
\label{eq:13-93}
\frac{\partial u_y}{\partial t} + 2\Omega_z u_x = 0,
\end{equation}
\begin{equation}
\label{eq:13-94}
\frac{\partial u_z}{\partial t} + \frac{c^2}{\rho}\frac{\partial}{\partial z}\delta\rho - \frac{\partial}{\partial z}\delta\mathfrak{B} - 2\Omega_y u_x = 0,
\end{equation}
\begin{equation}
\label{eq:13-95}
\frac{\partial}{\partial t}\delta\rho + \rho\frac{\partial u_z}{\partial z} = 0,
\end{equation}
and
\begin{equation}
\label{eq:13-96}
\frac{\partial^2}{\partial z^2}\delta\mathfrak{B} + 4\pi G\delta\rho = 0.
\end{equation}

Seeking solutions of these equations whose $(z,t)$ dependence is given by
\begin{equation}
\label{eq:13-97}
e^{i(kz+\sigma t)},
\end{equation}
we are led to the following system of linear homogeneous equations:
\begin{equation}
\label{eq:13-98}
\begin{vmatrix}
2i\Omega_z & \sigma & -2i\Omega_y & 0 & 0 \\
\sigma & -2i\Omega_z & 0 & 0 & 0 \\
0 & +2i\Omega_y & \sigma & c^2k/\rho & -k \\
0 & 0 & pk & \sigma & 0 \\
0 & 0 & 0 & 4\pi G & -k^2
\end{vmatrix}
\begin{pmatrix}
u_y \\ u_x \\ u_z \\ \delta\rho \\ \delta\mathfrak{B}
\end{pmatrix} = 0.
\end{equation}
The determinant of this system of equations must vanish; and on expanding the determinant and reducing, we obtain the characteristic equation
\begin{equation}
\label{eq:13-99}
\sigma^4 - (4\Omega^2+c^2k^2-4\pi G\rho)\sigma^2 + 4\Omega^2(c^2k^2-4\pi G\rho)\cos^2\!\vartheta = 0,
\end{equation}
where $\vartheta$ is the inclination of the direction of $\boldsymbol{\Omega}$ to the direction of wave propagation.

From equation~\eqref{eq:13-99} it follows that there are in general two modes in which a wave can be propagated in the medium. If $\sigma_1$ and $\sigma_2$ are the gyration frequencies of the two modes, then,
\begin{equation}
\label{eq:13-100}
\sigma_1^2+\sigma_2^2 = 4\Omega^2+c^2k^2-4\pi G\rho
\end{equation}
and
\begin{equation}
\label{eq:13-101}
\sigma_1^2\sigma_2^2 = 4\Omega^2(c^2k^2-4\pi G\rho)\cos^2\!\vartheta.
\end{equation}
It is also clear that both the roots $\sigma_1^2$ and $\sigma_2^2$ are real.

From equation~\eqref{eq:13-101} we can now conclude that if Jeans's condition,
\begin{equation}
\label{eq:13-102}
c^2k^2-4\pi G\rho < 0,
\end{equation}
is satisfied, then one of the two roots $\sigma_1^2$ or $\sigma_2^2$ must be negative; and, accordingly, one of the two modes must be unstable. An exception arises when the direction of wave propagation is at right angles to the direction of $\boldsymbol{\Omega}$; for, in this case, $\vartheta = \frac{1}{2}\pi$ and equation~\eqref{eq:13-99} gives
\begin{equation}
\label{eq:13-103}
\sigma^2 = 4\Omega^2+c^2k^2-4\pi G\rho;
\end{equation}
and the condition for the instability of these waves, namely,
\begin{equation}
\label{eq:13-104}
4\Omega^2+c^2k^2-4\pi G\rho < 0,
\end{equation}
cannot be met if $\Omega^2 > \pi G\rho$. We have thus shown: \textit{for waves propagated in a direction at right angles to the direction of $\boldsymbol{\Omega}$, gravitational instability cannot occur if $\Omega^2 > \pi G\rho$; for propagation in every other direction gravitational instability will occur if Jeans's condition, $c^2k^2-4\pi G\rho < 0$, is satisfied.}

In considering equation~\eqref{eq:13-99} generally, it is convenient to measure $\sigma$ in the unit $\sqrt{4\pi G\rho}$ sec$^{-1}$,
\begin{equation}
\label{eq:13-105}
k \text{ in the unit }\frac{1}{c}\sqrt{4\pi G\rho}\text{ cm}^{-1},
\end{equation}
and $\sigma/k$ in the unit $c$ cm/sec. In these units, Jeans's criterion is, for example, $k < 1$.

When $\sigma$ and $k$ are measured in the units specified in~\eqref{eq:13-105}, equation~\eqref{eq:13-99} takes the form
\begin{equation}
\label{eq:13-106}
\sigma^4 - (\Lambda^2+k^2-1)\sigma^2+\Lambda^2(k^2-1)\cos^2\!\vartheta = 0,
\end{equation}
where
\begin{equation}
\label{eq:13-107}
\Lambda = \Omega/\sqrt{\pi G\rho}
\end{equation}
is the non-dimensional combination in which $\Omega$ enters this problem. The roots of equation~\eqref{eq:13-106} are given by
\begin{equation}
\label{eq:13-108}
2\sigma^2 = \Lambda^2+k^2-1 \pm \sqrt{(\Lambda^2+k^2-1)^2-4\Lambda^2(k^2-1)\cos^2\!\vartheta}.
\end{equation}
The dispersion relations derived from this equation for $\Lambda^2 = 0.5$, 1.0, and 2.0 and $\vartheta = 0^\circ$, $45^\circ$, and $90^\circ$ are illustrated in Fig.~135. The manner in which the two modes cross one another and fill restricted parts of the $\sigma$, $k$-plane is curious and worth noting.

\subsection*{(b) \textit{The effect of a uniform magnetic field}}

When a uniform field of intensity $\mathbf{H}$ pervades the whole medium, the relevant perturbation equations are
\begin{equation}
\label{eq:13-109}
\rho\frac{\partial\mathbf{u}}{\partial t} = \frac{1}{4\pi}\operatorname{curl}\mathbf{h}\times\mathbf{H} - c^2\operatorname{grad}\delta\rho + \rho\operatorname{grad}\delta\mathfrak{B},
\end{equation}
\begin{equation}
\label{eq:13-110}
\frac{\partial\mathbf{h}}{\partial t} = \operatorname{curl}(\mathbf{u}\times\mathbf{H}),
\end{equation}
\begin{equation}
\label{eq:13-111}
\frac{\partial}{\partial t}\delta\rho = -\rho\operatorname{div}\mathbf{u}, \quad \text{and} \quad \nabla^2\delta\mathfrak{B} = -4\pi G\delta\rho.
\end{equation}
We shall seek solutions of equations~\eqref{eq:13-109}--\eqref{eq:13-111} which correspond to the propagation of waves in the $z$-direction. Then from the solenoidal character of $\mathbf{h}$ it follows that $h_z = 0$. And by further choosing the orientation of the coordinate system such that
\begin{equation}
\label{eq:13-112}
\mathbf{H} = (0, H_y, H_z),
\end{equation}

\figurePlaceholder{135}{The effect of rotation on the gravitational instability of an infinite homogeneous medium. The dispersion relations derived from equation~\eqref{eq:13-106} are illustrated for values of $\Lambda^2\,(= \Omega^2/\pi G\rho) = 0.5$, 1.0, and 2.0 in figures (a), (b), and (c), respectively. The curves are labelled by the angle which the direction of wave propagation makes with the direction $\boldsymbol{\Omega}$. The modes belonging to the same branch are drawn in the same way: either full-line or dashed. The abscissa gives the wave number in the unit $c^{-1}\sqrt{4\pi G\rho}$ cm$^{-1}$ while the ordinate gives $\sigma/|c|$ with $\sigma$ measured in the unit $\sqrt{4\pi G\rho}$ sec$^{-1}$. (Negative values of the ordinate correspond to instability.)}

we find that the equations which follow from equations~\eqref{eq:13-109}--\eqref{eq:13-111} break up into the two non-combining systems:
\begin{equation}
\label{eq:13-113}
\rho\frac{\partial u_x}{\partial t} = \frac{H_z}{4\pi}\frac{\partial h_x}{\partial z}, \quad \frac{\partial h_x}{\partial t} = H_z\frac{\partial u_x}{\partial z},
\end{equation}
\begin{equation}
\label{eq:13-114}
\rho\frac{\partial u_y}{\partial t} - \frac{H_z}{4\pi}\frac{\partial h_y}{\partial z} = 0,
\end{equation}
\begin{equation}
\label{eq:13-115}
\rho\frac{\partial u_z}{\partial t} - \frac{H_y}{4\pi}\frac{\partial h_y}{\partial z} + c^2\frac{\partial}{\partial z}\delta\rho - \rho\frac{\partial}{\partial z}\delta\mathfrak{B} = 0,
\end{equation}
\begin{equation}
\label{eq:13-116}
\frac{\partial h_y}{\partial t} + H_z\frac{\partial u_y}{\partial z} - H_y\frac{\partial u_z}{\partial z} = 0,
\end{equation}
\begin{equation}
\label{eq:13-117}
\frac{\partial}{\partial t}\delta\rho + \rho\frac{\partial u_z}{\partial z} = 0,
\end{equation}
\begin{equation}
\label{eq:13-118}
\frac{\partial^2}{\partial z^2}\delta\mathfrak{B} + 4\pi G\delta\rho = 0.
\end{equation}

Equations~\eqref{eq:13-113} can be combined to give
\begin{equation}
\label{eq:13-119}
\frac{\partial^2 h_x}{\partial t^2} = \frac{H_z^2}{4\pi\rho}\frac{\partial^2 h_x}{\partial z^2} \quad \text{and} \quad \frac{\partial^2 u_x}{\partial t^2} = \frac{H_z^2}{4\pi\rho}\frac{\partial^2 u_x}{\partial z^2}.
\end{equation}
These equations are the same as those leading to the ordinary hydromagnetic waves of Alfv\'en propagated with the velocity $H_z/\sqrt{4\pi\rho}$ (cf.\ Chapter~IV, \S\,39).

Turning next to equations~\eqref{eq:13-114}--\eqref{eq:13-118} and seeking solutions having a $(z,t)$-dependence given by~\eqref{eq:13-97}, we are led to the following system of linear homogeneous equations:
\begin{equation}
\label{eq:13-120}
\begin{vmatrix}
\rho\sigma & -kH_z/4\pi & 0 & 0 & 0 \\
0 & +kH_z/4\pi\rho & \rho\sigma & kc^2 & -k\rho \\
-H_y & 0 & kH_z & \sigma & 0 \\
0 & k\rho & 0 & \sigma & 0 \\
0 & 0 & 0 & 4\pi G & -k^2
\end{vmatrix}
= 0.
\end{equation}

The vanishing of the determinant of this system leads to the characteristic equation
\begin{equation}
\label{eq:13-121}
\sigma^4 - \frac{H^2k^2}{4\pi\rho}\sigma^2 + \frac{H_z^2k^2}{4\pi\rho}(c^2k^2-4\pi G\rho)\cos^2\!\vartheta = 0,
\end{equation}
where $\vartheta$ is now the inclination of the direction of $\mathbf{H}$ to the direction of wave propagation. From equation~\eqref{eq:13-121} it now follows that there are in general two modes of wave propagation besides the pure Alfv\'en modes described by equations~\eqref{eq:13-119}. If $\sigma_1$ and $\sigma_2$ are the gyration frequencies of the two modes described by equation~\eqref{eq:13-121}, then
\begin{equation}
\label{eq:13-122}
\sigma_1^2+\sigma_2^2 = \frac{H^2k^2}{4\pi\rho} + c^2k^2 - 4\pi G\rho
\end{equation}
and
\begin{equation}
\label{eq:13-123}
\sigma_1^2\sigma_2^2 = \frac{H_z^2k^2}{4\pi\rho}(c^2k^2-4\pi G\rho)\cos^2\!\vartheta.
\end{equation}
It is also clear that both the roots $\sigma_1^2$ and $\sigma_2^2$ are real.

From equation~\eqref{eq:13-123} we can now conclude that if Jeans's criterion~\eqref{eq:13-102} is satisfied, then one of the two roots $\sigma_1^2$ or $\sigma_2^2$ is negative; and accordingly, one of the two modes must be unstable. Thus, \textit{Jeans's condition for gravitational instability is unaffected by the presence of a uniform magnetic field.} The physical reason for this is evident for a mode in which the density waves are perpendicular to the lines of force; for, in that case the motion of the particles are parallel to the lines of force and are, therefore, unaffected by the magnetic field.

In considering equation~\eqref{eq:13-121} generally, it is convenient to measure $\sigma$ and $k$ in the units specified in~\eqref{eq:13-105}. Equation~\eqref{eq:13-121} then becomes
\begin{equation}
\label{eq:13-124}
\sigma^4 - \left(\frac{V_A^2}{c^2}k^2+k^2-1\right)\sigma^2 + \frac{V_A^2}{c^2}k^2(k^2-1)\cos^2\!\vartheta = 0,
\end{equation}
where
\begin{equation}
\label{eq:13-125}
V_A = H/\sqrt{4\pi\rho}
\end{equation}
is the Alfv\'en speed.

The dispersion relations derived from equation~\eqref{eq:13-124} for $V_A^2/c^2 = 0.5$, 1.0, and 2.0 and $\vartheta = 0^\circ$, $45^\circ$, and $90^\circ$ are illustrated in Fig.~136. The

\figurePlaceholder{136}{The effect of a uniform magnetic field on the gravitational instability of an infinite homogeneous medium. The dispersion relation derived from equation~\eqref{eq:13-124} is illustrated for values of $V_A^2/c^2\,(= H^2/(4\pi\rho c^2)) = 0.5$, 1.0, and 2.0 in figures (a), (b), and (c), respectively. The curves are labelled by the angle which the direction of wave propagation makes with the direction of the field. The modes belonging to the same branch are drawn in the same way: either full-line or dashed. The abscissa gives the wave number in the unit $c^{-1}\sqrt{4\pi G\rho}$ cm$^{-1}$ while the ordinate gives $\sigma/|c|$ with $\sigma$ measured in the unit $\sqrt{4\pi G\rho}$ sec$^{-1}$. (Negative values of the ordinate correspond to instability.)}

manner in which the two modes cross one another and fill restricted parts of the $\sigma$, $k$-plane is worth noting. The principal fact concerning the crossings is that the modes which, for $H \to 0$, are the pure Jeans and the pure Alfv\'en modes become, when $H \to \infty$, \textit{modified Alfv\'en} and \textit{modified Jeans} modes, respectively; the `modification' consists in the replacement of $H_z$ by $H$ in the expression for the Alfv\'en velocity in the former case and in the addition of a factor $\cos\vartheta$ to the Jeans velocity in the latter case. This crossing of the two modes is not relevant to the unstable modes which always occur for $k < 1$ (in the chosen units).

\subsection*{(c) \textit{The effect of the simultaneous presence of rotation and magnetic field}}

Finally, we shall consider the joint effects of a uniform rotation and a uniform magnetic field on the gravitational instability of an infinite homogeneous medium. In this general case, equation~\eqref{eq:13-109} is replaced by
\begin{equation}
\label{eq:13-126}
\rho\frac{\partial\mathbf{u}}{\partial t} = \frac{1}{4\pi}\operatorname{curl}\mathbf{h}\times\mathbf{H} + 2\rho\,\mathbf{u}\times\boldsymbol{\Omega} - c^2\operatorname{grad}\delta\rho + \rho\operatorname{grad}\delta\mathfrak{B};
\end{equation}
but the remaining equations~\eqref{eq:13-110} and~\eqref{eq:13-111} are unchanged.

Choosing the orientation of the coordinate axes such that
\begin{equation}
\label{eq:13-127}
\mathbf{H} = (0, H_y, H_z) \quad \text{and} \quad \boldsymbol{\Omega} = (\Omega_x, \Omega_y, \Omega_z),
\end{equation}
and seeking solutions which are independent of $x$ and $y$ and whose dependence on $z$ and $t$ is given by~\eqref{eq:13-97}, we are led to the following system of linear homogeneous equations:
\begin{equation}
\label{eq:13-128}
\begin{vmatrix}
\sigma & 0 & -kH_z & 0 & 0 & 0 \\
0 & \sigma & -kH_z & kH_y & 0 & 0 \\
-kH_z/4\pi\rho & 0 & \sigma & +2i\Omega_z & -2i\Omega_y & 0 \\
0 & +kH_z/4\pi\rho & +2i\Omega_z & \sigma & +2i\Omega_x & 0 \\
0 & 0 & 0 & \rho k & \sigma & 0 \\
0 & 0 & 0 & 0 & 4\pi G & -k^2
\end{vmatrix}
= 0.
\end{equation}

The vanishing of the determinant of this system leads to the characteristic equation
\begin{equation}
\label{eq:13-129}
\begin{split}
\sigma^6 &- \sigma^4\{(4\Omega^2+\Omega_1^2)+(\Omega_h^2+\Omega_3^2)\} + \\
&+ \sigma^2\{\Omega_1^2(4\Omega^2+\Omega_3^2)+\Omega_h^2(4\Omega^2+\Omega_3^2+\Omega_1^2+\Omega_3^2)\} + 4(\Omega_x\Omega_z - \Omega_y\Omega_x)^2 - \\
&\quad -\Omega_h^2\Omega_1^2\Omega_3^2 = 0,
\end{split}
\end{equation}
where
\begin{equation}
\label{eq:13-130}
\Omega_1^2 = \frac{H^2k^2}{4\pi\rho}, \quad \Omega_h^2 = \frac{H_z^2k^2}{4\pi\rho}, \quad \Omega^2 = |\boldsymbol{\Omega}|^2,
\end{equation}
and
\begin{equation*}
\Omega_3^2 = c^2k^2-4\pi G\rho.
\end{equation*}

From equation~\eqref{eq:13-129} it follows that there are in general three modes in which a wave can be propagated in the medium; and if $\sigma_1$, $\sigma_2$, and $\sigma_3$ denote the gyration frequencies of the three modes, then
\begin{equation}
\label{eq:13-131}
\sigma_1^2+\sigma_2^2+\sigma_3^2 = 4\Omega^2 + \frac{(H_y^2+H_z^2)k^2}{4\pi\rho} + c^2k^2 - 4\pi G\rho
\end{equation}
and
\begin{equation}
\label{eq:13-132}
\sigma_1^2\sigma_2^2\sigma_3^2 = \Omega_h^2\Omega_3^2 = \frac{H_z^2k^2}{4\pi\rho}(c^2k^2-4\pi G\rho).
\end{equation}
From equation~\eqref{eq:13-132} (or, more directly from equation~\eqref{eq:13-129}) it follows that if $\Omega_3^2$ is negative, then one of the roots of $\sigma^6$ must be negative; and this means that one of the three modes is unstable. The theorem stated at the beginning of the section follows.

\section*{BIBLIOGRAPHICAL NOTES}

\textit{For an exceptionally complete and penetrating analysis of the problems bearing on questions of stability as they occur in astrophysics, see:}

\begin{enumerate}
\item P.\ Ledoux, `Stellar stability', \textit{Handbuch der Physik}, \textbf{51}, 605--88 (1958), see \S\,17.
\end{enumerate}

\textit{Related problems on stellar variability and stellar pulsation are treated in:}

\begin{enumerate}
\setcounter{enumi}{1}
\item P.\ Ledoux and T.\ Walraven, `Variable stars', ibid.\ \textbf{363--604} (1958), see \S\S\,80 and 82.
\end{enumerate}

\S\,117. The scalar form of the virial theorem for an assembly of particles interacting according to some general law of force is well known in statistical mechanics; see, for example:

\begin{enumerate}
\setcounter{enumi}{2}
\item R.\ H.\ Fowler, \textit{Statistical Mechanics}, \S\,9.7, Cambridge, England, 1936.
\end{enumerate}

Essentially the same theorem has been known in celestial mechanics and classical dynamics for a very long time; it was discovered by Lagrange as a useful integral formula for the theory of the three-body problem; and it was generalized to $n$ bodies by Jacobi:

\begin{enumerate}
\setcounter{enumi}{3}
\item J.\ L.\ Lagrange, \textit{Oeuvres}, vi, 240; ix, 856, Paris, 1873.
\item C.\ G.\ J.\ Jacobi, \textit{Vorlesungen \"uber Dynamik}, pp.\ 22--23, Druck \& Verlag von Georg Reimer, Berlin, 1866.
\end{enumerate}

Lagrange's identity, as it is referred to in books on celestial mechanics, plays an important role in the modern developments of the subject:

\begin{enumerate}
\setcounter{enumi}{5}
\item C.\ L.\ Siegel, \textit{Vorlesungen \"uber Himmelsmechanik}, \S\,6, Springer-Verlag, Berlin, 1956.
\end{enumerate}

The usefulness of the virial theorem for stellar dynamics was recognized by

\begin{enumerate}
\setcounter{enumi}{6}
\item H.\ Poincar\'e, \textit{Le\c{c}ons sur les hypoth\`eses cosmogoniques}, \S\,74, Paris, 1911.
\item A.\ S.\ Eddington, `The kinetic energy of a star cluster', \textit{Mon.\ Not.\ Roy.\ Astr.\ Soc.}\ \textbf{76}, 525--8 (1916).
\end{enumerate}

For an account of these latter applications see:

\begin{enumerate}
\setcounter{enumi}{8}
\item S.\ Chandrasekhar, \textit{Principles of Stellar Dynamics}, chap.\ v, University of Chicago Press, 1942.
\end{enumerate}

An extension of the virial theorem to include frictional terms proportional to the velocity was given by Milne:

\begin{enumerate}
\setcounter{enumi}{9}
\item E.\ A.\ Milne, `An extension of the theorem of the virial', \textit{Phil.\ Mag.}, Ser.\ 6, \textbf{50}, 409--14 (1925).
\end{enumerate}

The suggestion that one might formulate the virial theorem in tensor form seems to have been made, first, by Rayleigh:

\begin{enumerate}
\setcounter{enumi}{10}
\item Lord Rayleigh, `On a theorem analogous to the virial theorem', \textit{Scientific Papers}, iv, 491--3, Cambridge, England, 1903.
\end{enumerate}

More recently, this aspect has been revived by:

\begin{enumerate}
\setcounter{enumi}{11}
\item E.\ N.\ Parker, `Tensor virial equations', \textit{Physical Rev.}\ \textbf{96}, 1686--9 (1954).
\end{enumerate}

The extension of the virial theorem to hydromagnetics was given by:

\begin{enumerate}
\setcounter{enumi}{12}
\item S.\ Chandrasekhar and E.\ Fermi, `Problems of gravitational stability in the presence of a magnetic field', \textit{Astrophys.\ J.}\ \textbf{118}, 116--41 (1953).

\item[13a.] S.\ Chandrasekhar, `The virial theorem in hydromagnetics', \textit{Journal of Mathematical Analysis and Applications}, \textbf{1}, 240--53 (1960).
\end{enumerate}

See also:

\begin{enumerate}
\setcounter{enumi}{13}
\item S.\ Chandrasekhar, `Problems of stability in hydrodynamics and hydromagnetics', \textit{Mon.\ Not.\ Roy.\ Astr.\ Soc.}\ \textbf{113}, 667--78 (1953).
\end{enumerate}

In reference 13 the scalar form of the virial theorem (equation~\eqref{eq:13-29} in the text) was derived; and its consequences for gravitational equilibrium (in particular the inequality, $\mathfrak{M} < |\mathfrak{W}|$) were discussed. The general treatment of the virial theorem given in the text follows reference 13a.

\S\,118. The first use of the virial theorem to obtain an integral formula for $\sigma^2$ is contained in:

\begin{enumerate}
\setcounter{enumi}{14}
\item P.\ Ledoux, `On the radial pulsation of gaseous stars', \textit{Astrophys.\ J.}\ \textbf{102}, 143--53 (1945).
\end{enumerate}

In this paper Ledoux derives the formula $\sigma^2 = -(3\gamma-4)\mathfrak{W}/I$. In this particular instance, Ledoux was able to show that the same formula follows, also, from a different integral relation provided by a variational treatment of the underlying characteristic value problem for $\sigma^2$. For this latter treatment see:

\begin{enumerate}
\setcounter{enumi}{15}
\item P.\ Ledoux and C.\ L.\ Pekeris, `Radial pulsations of stars', ibid.\ \textbf{94}, 124--35 (1941).
\end{enumerate}

The generalization of Ledoux's method to allow for prevalent magnetic fields is contained in:

\begin{enumerate}
\setcounter{enumi}{16}
\item S.\ Chandrasekhar and D.\ Nelson Limber, `On the pulsation of a star in which there is a prevalent magnetic field', ibid.\ \textbf{119}, 10--13 (1954).
\end{enumerate}

In this paper the formula $\sigma^2 = -(3\gamma-4)(\mathfrak{W}+\mathfrak{M})/I$ is derived. The inconsistencies in the assumption of a trial function for $\boldsymbol{\xi}$ which corresponds to a uniform radial expansion, are avoided by the substitution $\boldsymbol{\xi} = \mathsf{X}\mathbf{x}$ considered in the text.

See also:

\begin{enumerate}
\setcounter{enumi}{17}
\item E.\ N.\ Parker, `The gross dynamics of a hydromagnetic gas cloud', \textit{Astrophys.\ J.\ Suppl.}, Ser.\ 3, \textbf{51}--76 (1957).
\end{enumerate}

\S\,119. As stated in the text, the gravitational instability of an infinite homogeneous medium was discovered by Jeans. His first statement of the theorem occurs in:

\begin{enumerate}
\setcounter{enumi}{18}
\item J.\ H.\ Jeans, `The stability of a spherical nebula', \textit{Phil.\ Trans.\ Roy.\ Soc.\ (London)}, \textbf{199}, 1--53 (1902), see particularly, \S\,46.
\end{enumerate}

For a later account by the same author see:

\begin{enumerate}
\setcounter{enumi}{19}
\item J.\ H.\ Jeans, \textit{Astronomy and Cosmogony}, p.\ 313, Cambridge, 1929.
\end{enumerate}

The effect of turbulence on Jeans's criterion is considered by:

\begin{enumerate}
\setcounter{enumi}{20}
\item S.\ Chandrasekhar, `The gravitational instability of an infinite homogeneous turbulent medium', \textit{Proc.\ Roy.\ Soc.\ (London)} A, \textbf{210}, 26--29 (1951).
\end{enumerate}

\S\,120. An excellent and an original review of the effects of rotation and magnetic field on gravitational instability has been written by:

\begin{enumerate}
\setcounter{enumi}{21}
\item A.\ G.\ Pacholczyk and J.\ S.\ Stodkiewicz, `On the gravitational instability of some magnetohydrodynamical systems of astrophysical interest', \textit{Acta Astronomica} (Polska Akademia Nauk) \textbf{10}, 1--29 (1960).
\end{enumerate}

\S\,120\,(a). The effect of rotation on Jeans's criterion was considered by:

\begin{enumerate}
\setcounter{enumi}{22}
\item S.\ Chandrasekhar, `The gravitational instability of an infinite homogeneous medium when a Coriolis acceleration is acting', \textit{Vistas in Astronomy I}, 344--7, edited by A.\ Beer, Pergamon Press, 1955.
\end{enumerate}

The illustration of the dispersion relations in the manner of Fig.~135 follows a suggestion of Pacholczyk and Stodkiewicz (reference 22). In their account these latter authors include the effects of viscosity and thermal conductivity.

The extension of Jeans's criterion to nonuniform rotation is considered by:

\begin{enumerate}
\setcounter{enumi}{23}
\item E.\ Schatzman and N.\ Bel, `Instabilit\'e d'une masse fluide \'etendue', \textit{Comptes rendus des s\'eances de l'Acad\'emie des Sciences}, \textbf{241}, 20--22 (1955).
\item N.\ Bel, `Instabilit\'e d'une masse fluide \'etendue', ibid.\ 163--4 (1955).
\end{enumerate}

See also:

\begin{enumerate}
\setcounter{enumi}{25}
\item P.\ Ledoux, `Sur la stabilit\'e gravitationnelle d'une n\'ebuleuse isotherme', \textit{Ann.\ d'Astrophys.}\ \textbf{14}, 438--47 (1951).
\item W.\ Fricke, `On the gravitational instability in a rotating isothermal medium', \textit{Astrophys.\ J.}\ \textbf{120}, 356--9 (1954).
\end{enumerate}

\S\,120\,(b). The effect of a uniform magnetic field on Jeans's criterion was considered by Chandrasekhar and Fermi (reference 13). Again, the illustration of the dispersion relations in the manner of Fig.~136 follows a suggestion of Pacholczyk and Stodkiewicz (reference 22). These latter authors have discussed, at the same time, the effects of finite electrical conductivity.

\S\,120\,(c). The joint effects of rotation and magnetic field are considered by:

\begin{enumerate}
\setcounter{enumi}{27}
\item S.\ Chandrasekhar, `The gravitational instability of an infinite homogeneous medium when Coriolis force is acting and a magnetic field is present', \textit{Astrophys.\ J.}\ \textbf{119}, 7--9 (1954).
\end{enumerate}
